\documentclass[a4paper,10.5pt]{article}
\usepackage[T1]{fontenc}\usepackage[utf8]{inputenc}
\usepackage[top=2.2cm,bottom=2.2cm,left=2.35cm,right=2.35cm]{geometry}
\usepackage{hyperref}\usepackage{xcolor}\definecolor{accent}{HTML}{1F4E78}\hypersetup{colorlinks=true,urlcolor=accent}
\pagestyle{empty}\setlength{\parindent}{0pt}\setlength{\parskip}{7pt}
\begin{document}
\textbf{Prabhavit Gupta}\\
Venice, Italy \textbar{} +39 351 610 7807 \textbar{} prabhavitg@gmail.com\\[10pt]
Dear Konfetti Team,

I am writing to apply for the Backend Developer (Laravel / AWS) role in Berlin. A lean team with real ownership of technical direction, backed by Laravel and AI tooling used day to day, is exactly the kind of setup I have been working in for the past three years.

At FlowSpace, I have been the lead developer on a live workspace marketplace built on Laravel and Next.js, owning backend features end to end: I drove the migration off a legacy stack onto Laravel, mapping the old data model to the new one and cutting over while the product stayed live. Claude Code is a core part of how I build there, not just a helper on the side: I used it to build and still own an AI listing bot that turns a provider brochure into a drafted listing, and support AI agents that answer customer queries from live platform data. I also shipped the OfficeRnD and Nexudus API integrations for automated bookings, so I am used to working across service boundaries and external integrations, not just a single service.

At Urlink, I worked as a freelance developer and project manager, building the platform and its logic end to end with AI-assisted development and turning client requirements directly into shipped automations, in a small team where I owned the decisions rather than waiting on a ticket queue. That is close to the maker mindset and short idea-to-production path Konfetti describes.

I want to be upfront on AWS: my production experience is with Laravel and Supabase/PostgreSQL rather than AWS services like Lambda, SQS or ECS/Fargate. I have used AWS at a working-knowledge level, and I am ready to go deep on it quickly given how central Laravel and AI-assisted development already are to how I work.

For transparency: I hold Italian EU long-term resident status. I am not currently authorised to work in Germany and would need the applicable employer-side permit or sponsorship process. I am ready to relocate to Berlin for an office-first role and complete that process.

I would welcome the opportunity to discuss how I could contribute to the konfetti platform.

Regards,\\
Prabhavit Gupta
\end{document}
